\subsubsection{Terme de couture dans la concaténation \(108\to1080\)
}
\label{vi:subsubsec:costura-108-1080}

Le prolongement de la lecture chronologique doit conserver le prédécesseur
effectif de chaque événement. Soient
\(d^{(b)}=(d^{(b)}_1,\ldots,d^{(b)}_{108})\), \(0\le b<10\), dix
mots consécutifs, et définissons

\[
 a_b=d^{(b)}_1,\qquad q_b=d^{(b)}_{108},\qquad q_{-1}=q_9.
\]
Notons \(T_z^{\sqcup}\) la somme des dix lecteurs évalués
séparément, en utilisant \(q_b\) comme prédécesseur cyclique de \(a_b\), et
\(T_z^{\Vert}\) le lecteur du mot concaténé, qui utilise
\(q_{b-1}\) comme prédécesseur effectif de \(a_b\).


\begin{lemma}[Correction exacte de couture
]
\label{vi:lem:costura-108-1080}
Avec \(\Delta T_z=T_z^{\Vert}-T_z^{\sqcup}\),

\begin{equation}
 \Delta T_z=
 \sum_{b=0}^{9}\mathbf1_{\{a_b=9\}}
 \bigl(j_2(q_{b-1})-j_2(q_b)\bigr).
 \label{vi:eq:delta-costura-tz}
\end{equation}
Si les couronnes sont situées en \(t\equiv0\pmod{27}\), alors

\begin{equation}
 \Delta C_z=0,
 \qquad
 \Delta k_{\Delta_4}=\Delta T_z.
 \label{vi:eq:delta-costura-k}
\end{equation}
\end{lemma}

\begin{proof}
Hors des dix débuts de bloc, chaque événement conserve le même
prédécesseur dans les deux évaluations. Si \(a_b\ne9\), l’indicatrice qui définit
\(z_t\) annule également l’éventuelle différence de prédécesseur. Si \(a_b=9\),
le terme séparé est \(j_2(q_b)\) et le terme concaténé est
\(j_2(q_{b-1})\) ; leur différence est le terme de
\eqref{vi:eq:delta-costura-tz}. La somme des dix blocs donne la
première identité. Les nouveaux débuts occupent les positions
\(108b+1\equiv1\pmod{27}\), jamais une couronne. Par conséquent, la couture ne
modifie pas \(C_z\), et la définition
\(k_{\Delta_4}=T_z-2C_z\) donne
\eqref{vi:eq:delta-costura-k}.

\end{proof}

La correction atteint réellement la valeur \(-2\). Il suffit de prendre un unique
début \(a_b=9\), avec \(q_{b-1}=1\), \(q_b=2\), et d’exiger
\(a_r\ne9\) aux autres débuts. Alors
\(j_2(q_{b-1})-j_2(q_b)=-1-1=-2\). Ce mot est un contrôle du lecteur
de couture ; il n’est promu ni en route survivante, ni en particule, ni en signature
physique.


La longueur \(1080\) répète dix stations de phase, non dix copies
d’une profondeur invariable. Le transport correct est la composition
ordonnée

\[
 \mathcal H_9\circ\mathcal H_8\circ\cdots\circ\mathcal H_0.
\]
Seule une preuve supplémentaire de stationnarité permettrait de la remplacer par
\(\mathcal H_{108}^{10}\). La formule de couture conserve précisément
la mémoire que cette substitution pourrait effacer.

